\section{Redundant Relations in the HQC Equational Theory}
%
\label{App:Redundancies}

Two presentations $( \Sigma, E )$ and $( \Gamma, Q )$ are equivalent if $P( \Sigma ) / E \cong P( \Gamma ) / Q$ as prop categories.
It can be shown that $( \Sigma, E )$ and $( \Gamma, Q )$ are equivalent if and only if $( \Gamma, Q )$ can be obtained from $( \Sigma, E )$ via a sequence of Tietze transformations~(see \cite{Clerc2015} for a similar case).
The transformations are as follow.
\begin{itemize}
\item \textbf{Removing a Generator}.
      Let $f \in \Sigma$, $\Gamma = \Sigma \setminus \{ f \}$, and $g \in P( \Gamma )( \dom( f ), \cod( f ) )$.
      If $f \approx_E g$, then $( \Sigma, E )$ is equivalent to $( \Gamma, Q )$ where $Q$ is obtained from $E$ by replacing every instance of $f$ by $g$.
\item \textbf{Removing a Relation}.
      Let $( f, g ) \in E$ and $Q = E \setminus \{ ( f, g ) \}$.
      If $f \approx_{Q} g$, then $( \Sigma, E )$ is equivalent to $( \Sigma, Q )$.
\end{itemize}
These transformations can then be used to prove that specific generators and relations in $( \Sigma_{\CatHQC}, E )$ are redundant.

Using these transformations, this section derives several minimal generating sets for $\CatHQC / E$.
It is well-known that all single-qubit circuits can be generated (up-to global phase) using any two of the three generators $\{ X( \theta ), Z( \theta ), H \}$.
This corresponds to single-qubit hierarchical circuits generated from
$\{
        \scalebox{0.65}{\circimg{tikz-cache/dagger__term1.pdf}},
        \scalebox{0.65}{\circimg{tikz-cache/dagger__term3.pdf}},
        \scalebox{0.65}{\circimg{tikz-cache/dagger__term5.pdf}}
\}$.
Regardless of the choice of generators, it is then possible to find circuits $U$ and $V$ such that $\sem{U}_H = Z$ and $\sem{V}_H = X$.
The circuit $U \odot V$ can then act as a controlled-not gate, and all multi-qubit circuits can be generated (up-to global phase).
Finally, the global phase gate can be used to generate all multi-qubit circuits exactly, as described in~\cref{Prop:Full}.
These minimal gate sets can be obtained via Tietze transformations.

First, it will be shown that $g_1 = \scalebox{0.65}{\circimg{tikz-cache/dagger__term4.pdf}}$ can be eliminated from $( \Sigma_{\CatHQC}, E )$.
First, note that $g_1$ appears as the left-hand side of \HqcTargMDef, and does not appear on the right-hand side of \HqcTargMDef.
This means that $g_1$ and \HqcTargMDef can be removed from $( \Sigma_{\CatHQC}, E )$ via a Tietze transformation.
Moreover, since $g_1$ does not appear in any other relations, then the relations in $E$ need not be modified.

Next, it will be shown that $g_2 = \scalebox{0.65}{\circimg{tikz-cache/dagger__term2.pdf}}$ can be eliminated from $( \Sigma_{\CatHQC}, E )$.
As in the case of $g_1$, there exists a relation \HqcOPhaseDef which can be used to eliminate $g_2$ via a Tietze transformation.
However, $g_2$ also appears in \HqcCtrlPair, and must therefore be modified.
To simplify the process, we will manually replace \HqcCtrlPair before removing $g_2$.
Observe that the following derivation holds for each $U \in \CatHQC( n, n )$, since $H$ is self-adjoint.
\begin{equation*}
    \scalebox{0.65}{\circimg{tikz-cache/hqc__eqn15_lhs.pdf}}
    \overunderset{\HQCCtrlCong}{\HqcOPhaseDef}{=}
    \scalebox{0.65}{\circimg{tikz-cache/reductions__eqn_15_mid.pdf}}
    \overset{\HqcCtrlDag}{=}
    \scalebox{0.65}{\circimg{tikz-cache/hqc__eqn15_lhs_alt.pdf}}
\end{equation*}
Using the derivation, it is then possible to introduce a new relation \text{E'(13)}, which can then be used to eliminate \HqcCtrlPair.
\begin{equation*}
    \scalebox{0.65}{\circimg{tikz-cache/hqc__eqn15_lhs_alt.pdf}}
    \overset{\text{E'(13)}}{=}
    \scalebox{0.65}{\circimg{tikz-cache/hqc__eqn15_rhs.pdf}}
\end{equation*}
Since $g_2$ does not appear in any instance of \text{E'(13)}, then $g_2$ can be removed using \HqcOPhaseDef without any further modifications to the relation set.
This yields a new presentation, $( \Gamma, Q )$ where $\Gamma$ and $Q$ are defined as follows.
\begin{align*}
    \Gamma &= \{
        \scalebox{0.65}{\circimg{tikz-cache/dagger__term1.pdf}},
        \scalebox{0.65}{\circimg{tikz-cache/dagger__term3.pdf}},
        \scalebox{0.65}{\circimg{tikz-cache/dagger__term5.pdf}},
        \;\scalebox{0.65}{\circimg{tikz-cache/dagger__term8.pdf}}
    \}
    &
    Q &= E \setminus \{
        \HqcTargMDef, \HqcOPhaseDef, \HqcCtrlPair
    \} \cup \{
        \text{E'(13)}
    \}
\end{align*}

It remains to be shown that $\Gamma$ can be reduced to the global phase gate with only two of the three single-qubit gates.
Before proceeding, it will be shown that \HqcHDef is redundant, and can be removed from the presentation.
In principle, \HqcHDef can be derived via several applications of \HqcEuler, though this process is tedious.
Instead, it should be noted that in the proof of completeness, \HqcHDef is used precisely once, to prove that the Hadamard gate is self-inverse.
Then, it suffices to prove that the Hadamard gate is self-inverse without the use of \HqcHDef.
It then follows from completeness, together with the soundness of $E$, that \HqcHDef is derivable from $Q \setminus \{ \HqcHDef \}$.
As a first step, \HqcEuler is used to derive an alternative definition for the Hadamard gate.
\begin{equation}
    \scalebox{0.65}{\circimg{tikz-cache/reductions__hdef_0.pdf}}
    \overset{\HqcPowZero}{=}
    \scalebox{0.65}{\circimg{tikz-cache/reductions__hdef_1.pdf}}
    \overset{\HqcPowZero}{=}
    \scalebox{0.65}{\circimg{tikz-cache/reductions__hdef_2.pdf}}
    \overset{\HqcEuler}{=}
    \scalebox{0.65}{\circimg{tikz-cache/reductions__hdef_3.pdf}}
    \label{Eqn:HadDef}
\end{equation}
This alternative definition can then be used to prove that the Hadamard gate is self-inverse.
\begin{align*}
    \scalebox{0.65}{\circimg{tikz-cache/reductions__hinv_1.pdf}}
    \overset{\text{\cref{Eqn:HadDef}}}{=}&
    \scalebox{0.65}{\circimg{tikz-cache/reductions__hinv_2.pdf}}
    \overset{\HqcEuler}{=}
    \scalebox{0.65}{\circimg{tikz-cache/reductions__hinv_3.pdf}}
    \overset{\HqcPowZero}{=}
    \scalebox{0.65}{\circimg{tikz-cache/reductions__hinv_4.pdf}}
    \overset{\HqcPowZero}{=}
    \scalebox{0.65}{\circimg{tikz-cache/reductions__hinv_5.pdf}}
    \\
    \overset{\HqcPowAdd}{=}&
    \scalebox{0.65}{\circimg{tikz-cache/reductions__hinv_6.pdf}}
    \overset{\HqcPowAdd}{=}
    \scalebox{0.65}{\circimg{tikz-cache/reductions__hinv_7.pdf}}
    \overset{\HqcPOrd}{=}
    \scalebox{0.65}{\circimg{tikz-cache/reductions__hinv_8.pdf}}
    \overset{\HqcFreePhase}{=}
    \scalebox{0.65}{\circimg{tikz-cache/reductions__hinv_9.pdf}}
    \overunderset{\HQCCtrlCong}{\HqcPOrd}{=}
    \scalebox{0.65}{\circimg{tikz-cache/reductions__hinv_10.pdf}}
    \overset{\HqcCtrlId}{=}
    \scalebox{0.65}{\circimg{tikz-cache/reductions__hinv_11.pdf}}
\end{align*}
In conclusion, $( \Gamma, Q )$ is equivalent to $( \Gamma, R )$ where $R = Q \setminus \{ \HqcHDef \}$.
It remains to be shown that $\Gamma$ can be further reduced.

\begin{figure}[t]
    \begin{framed}
    \textbf{Collection 1: Ordinary Circuit Relations}
    \vspace{-0.5em}
    \begin{align*}
        \scalebox{0.65}{\circimg{tikz-cache/proofs__decode_q3_lhs.pdf}}
        &\overset{\text{R'(5)}}{=}
        \scalebox{0.65}{\circimg{tikz-cache/hqc__eqn5_rhs_alt.pdf}}
        & % -------------------------------------------------------
        \scalebox{0.65}{\circimg{tikz-cache/hqc__euler_lhs_alt.pdf}}
        &\overset{\text{R'(6)}}{=}
        \scalebox{0.65}{\circimg{tikz-cache/hqc__euler_rhs_alt.pdf}}
        & % -------------------------------------------------------
        \scalebox{0.65}{\circimg{tikz-cache/hqc__eqn7_lhs.pdf}}
        &\overset{\HqcPOrd}{=}
        \scalebox{0.65}{\circimg{tikz-cache/cqc__eqn1_rhs.pdf}}
    \end{align*}
    
    \textbf{Collection 2: Generalized Control Relations}
    \vspace{-0.5em}
    \begin{align*}
        \scalebox{0.65}{\circimg{tikz-cache/hqc__eqn1_lhs.pdf}}
        &\overset{\HqcCtrlId}{=}
        \scalebox{0.65}{\circimg{tikz-cache/hqc__eqn1_rhs.pdf}}
        & % -------------------------------------------------------
        \scalebox{0.65}{\circimg{tikz-cache/hqc__eqn2_lhs.pdf}}
        &\overset{\HqcCtrlComp}{=}
        \scalebox{0.65}{\circimg{tikz-cache/hqc__eqn2_rhs.pdf}}
        & % -------------------------------------------------------
        \scalebox{0.65}{\circimg{tikz-cache/hqc__eqn3_lhs.pdf}}
        &\overset{\HqcCtrlSwap}{=}
        \scalebox{0.65}{\circimg{tikz-cache/hqc__eqn3_rhs.pdf}}
        & % -------------------------------------------------------
        \scalebox{0.65}{\circimg{tikz-cache/hqc__eqn18_lhs.pdf}}
        &\overset{\HqcFreePhase}{=}
        \scalebox{0.65}{\circimg{tikz-cache/hqc__eqn18_rhs.pdf}}
    \end{align*}
    \vspace{-1.5em}
    \begin{align*}
        \scalebox{0.65}{\circimg{tikz-cache/hqc__eqn14_lhs.pdf}}
        &\overset{\HqcCtrlDag}{=}
        \scalebox{0.65}{\circimg{tikz-cache/hqc__eqn14_rhs.pdf}}
        & % -------------------------------------------------------
        \scalebox{0.65}{\circimg{tikz-cache/hqc__eqn15_lhs_alt2.pdf}}
        &\overset{\HqcCtrlPair}{=}
        \scalebox{0.65}{\circimg{tikz-cache/hqc__eqn15_rhs.pdf}}
    \end{align*}
    \vspace{-1.5em}
    \begin{align*}
        \scalebox{0.65}{\circimg{tikz-cache/hqc__eqn22_lhs.pdf}}
        &\overset{\HqcCtrlExpand}{=}
        \scalebox{0.65}{\circimg{tikz-cache/hqc__eqn22_rhs.pdf}}
        & % -------------------------------------------------------
        \scalebox{0.65}{\circimg{tikz-cache/hqc__eqn23_lhs.pdf}}
        &\overset{\HqcCtrlMerge}{=}
        \scalebox{0.65}{\circimg{tikz-cache/hqc__eqn23_rhs.pdf}}
    \end{align*}
    
    \textbf{Collection 3: Generalized Exponential Relations}
    \vspace{-1em}
    \begin{align*}
        \scalebox{0.65}{\circimg{tikz-cache/hqc__eqn9_lhs.pdf}}
        &\overset{\HqcPowZero}{=}
        \scalebox{0.65}{\circimg{tikz-cache/hqc__eqn9_rhs.pdf}}
        & % -------------------------------------------------------
        \scalebox{0.65}{\circimg{tikz-cache/hqc__eqn4_lhs.pdf}}
        &\overset{\HqcPowId}{=}
        \scalebox{0.65}{\circimg{tikz-cache/hqc__eqn4_rhs.pdf}}
        & % -------------------------------------------------------
        \scalebox{0.65}{\circimg{tikz-cache/hqc__eqn5_lhs.pdf}}
        &\overset{\HqcPowAdd}{=}
        \scalebox{0.65}{\circimg{tikz-cache/hqc__eqn5_rhs.pdf}}
        & % -------------------------------------------------------
        \scalebox{0.65}{\circimg{tikz-cache/hqc__eqn24_lhs.pdf}}
        &\overset{\HqcPowIdem}{=}
        \scalebox{0.65}{\circimg{tikz-cache/hqc__eqn24_rhs.pdf}}
    \end{align*}
    \vspace{-2.5em}
    \begin{align*}
        \scalebox{0.65}{\circimg{tikz-cache/hqc__eqn19_lhs.pdf}}
        &\overset{\HqcPowDag}{=}
        \scalebox{0.65}{\circimg{tikz-cache/hqc__eqn19_rhs.pdf}}
        & % -------------------------------------------------------
        \scalebox{0.65}{\circimg{tikz-cache/hqc__eqn21_lhs.pdf}}
        &\overset{\HqcPowExpand}{=}
        \scalebox{0.65}{\circimg{tikz-cache/hqc__eqn21_rhs.pdf}}
    \end{align*}
    
    \textbf{Collection 4: Congruence Relations (Inductively Defined)}
    \vspace{-1em}
    \begin{equation*}
        \scalebox{0.65}{\circimg{tikz-cache/hqc__equiv_lhs.pdf}}
        \approx_E
        \scalebox{0.65}{\circimg{tikz-cache/hqc__equiv_rhs.pdf}}
        %
        \qquad\Rightarrow\qquad
        %
        \scalebox{0.65}{\circimg{tikz-cache/hqc__eqn20_lhs.pdf}}
        \overset{\HqcPowCong}{=}
        \scalebox{0.65}{\circimg{tikz-cache/hqc__eqn20_rhs.pdf}}
        %
        \qquad\text{and}\qquad
        %
        \scalebox{0.65}{\circimg{tikz-cache/hqc__eqn13_lhs.pdf}}
        \overset{\HQCCtrlCong}{=}
        \scalebox{0.65}{\circimg{tikz-cache/hqc__eqn13_rhs.pdf}}
    \end{equation*}
    
    \vspace{-1.5em}
\end{framed}

    \caption{An alternative equational theory for $\CatHQC$, using only the generators in $\Gamma_X$.}
    \label{Fig:RXRelations}
\end{figure}

\vspace{1em}
\noindent
\textbf{Eliminating the Pauli-X Gate}.
%
Let $g_X = \scalebox{0.65}{\circimg{tikz-cache/dagger__term3.pdf}}$.
Since $g_X$ appears as the left-hand side of \HqcTargDef, and does not appear on the right-hand side of \HqcTargDef, then $g_X$ and \HqcTargDef can be removed via a Tietze transformation. 
Unfortunately, $g_X$ still appears in relations $\{ \HqcSwapDef, \HqcEuler, \text{E'(13)} \}$, meaning that some care must be taken before $g_X$ can be fully removed from $( \Gamma, R )$.
To simplify the process, we will manually replace $\{ \HqcSwapDef, \HqcEuler, \text{E'(13)} \}$ before removing $g_X$.
This is possible since $H$ is self-adjoint.
To replace \HqcSwapDef, the following derivation suffices.
\begin{align*}
    \scalebox{0.65}{\circimg{tikz-cache/proofs__decode_q3_mid.pdf}}
    \overset{\HqcCtrlSwap}{=}&
    \scalebox{0.65}{\circimg{tikz-cache/proofs__decode_q3_rhs.pdf}}
    \overunderset{\HQCCtrlCong}{\HqcTargDef}{=}
    \scalebox{0.65}{\circimg{tikz-cache/reductions__swapdef_1.pdf}}
    \\
    \overset{\HqcCtrlDag}{=}&
    \scalebox{0.65}{\circimg{tikz-cache/reductions__swapdef_2.pdf}}
    \overunderset{\HQCCtrlCong}{\HqcTargDef}{=}
    \scalebox{0.65}{\circimg{tikz-cache/reductions__swapdef_3.pdf}}
    \\
    \overset{\HqcCtrlDag}{=}&
    \scalebox{0.65}{\circimg{tikz-cache/reductions__swapdef_4.pdf}}
    \overunderset{\HQCCtrlCong}{\HqcTargDef}{=}
    \scalebox{0.65}{\circimg{tikz-cache/reductions__swapdef_5.pdf}}
    \\
    \overset{\HqcCtrlDag}{=}&
    \scalebox{0.65}{\circimg{tikz-cache/hqc__eqn5_rhs_alt.pdf}}
\end{align*}
To replace \HqcEuler, the following derivations suffice.
\begin{align*}
    \scalebox{0.65}{\circimg{tikz-cache/hqc__euler_lhs.pdf}}
    &\overunderset{\HqcPowCong}{\HqcTargDef}{=}
    \scalebox{0.65}{\circimg{tikz-cache/reductions__euler_lhs_mid.pdf}}
    \overset{\HqcPowDag}{=}
    \scalebox{0.65}{\circimg{tikz-cache/hqc__euler_lhs_alt.pdf}}
    \\
    \scalebox{0.65}{\circimg{tikz-cache/hqc__euler_rhs.pdf}}
    &\overunderset{\HqcPowCong}{\HqcTargDef}{=}
    \scalebox{0.65}{\circimg{tikz-cache/reductions__euler_rhs_mid.pdf}}
    \overset{\HqcPowDag}{=}
    \scalebox{0.65}{\circimg{tikz-cache/hqc__euler_rhs_alt.pdf}}
\end{align*}
To replace \text{E'(13)}, the following derivation suffices.
\begin{align*}
    \scalebox{0.65}{\circimg{tikz-cache/hqc__eqn15_lhs_alt.pdf}}
    \overset{\HqcTargDef}{=}
    \scalebox{0.65}{\circimg{tikz-cache/reductions__eqn_15_mid_alt.pdf}}
    \overset{\HqcTargDef}{=}
    \scalebox{0.65}{\circimg{tikz-cache/hqc__eqn15_lhs_alt2.pdf}}
\end{align*}
Using these derivations, relations $\{ \text{R'(5)}, \text{R'(6)}, \text{R'(13)} \}$ can be introduced (see~\cref{Fig:RXRelations}), which can ten be used to eliminate $\{ \HqcSwapDef, \HqcEuler, \text{E'(13)} \}$.
Since $g_X$ does not appear in any instance of \text{E'(5)} or \text{E'(6)}, then $g_X$ can be removed using \HqcTargDef without any further modifications to the relation set.
This yields a new presentation, $( \Gamma_X, R_X )$ where $\Gamma_X$ and $R_X$ are defined as follows.
\begin{align*}
    \Gamma &= \{
        \scalebox{0.65}{\circimg{tikz-cache/dagger__term1.pdf}},
        \scalebox{0.65}{\circimg{tikz-cache/dagger__term5.pdf}},
        \;\scalebox{0.65}{\circimg{tikz-cache/dagger__term8.pdf}}
    \}
    &
    R_X &= R \setminus \{
        \HqcTargDef,
        \HqcSwapDef,
        \HqcEuler
    \} \cup \{
        \text{E'(5)},
        \text{E'(6)}
    \}
\end{align*}
The relations in $Q_X$ are depicted in \cref{Fig:RXRelations}.

\vspace{1em}
\noindent
\textbf{Eliminating the Pauli-Z Gate}.
%
Let $g_Z = \scalebox{0.65}{\circimg{tikz-cache/dagger__term1.pdf}}$.
A-priori, there does not exist a relation to eliminate $g_Z$.
However, the following derivation can be used to introduce such a relation, via a Tietze transformation.
\begin{align*}
    \scalebox{0.65}{\circimg{tikz-cache/reductions__zdef_1.pdf}}
    \overset{\HqcTargDef}{=}
    \scalebox{0.65}{\circimg{tikz-cache/reductions__zdef_2.pdf}}
    \approx_E
    \scalebox{0.65}{\circimg{tikz-cache/reductions__zdef_3.pdf}}
    \approx_E
    \scalebox{0.65}{\circimg{tikz-cache/reductions__zdef_4.pdf}}
\end{align*}
Since the generalized control relations in $( \Sigma_{\CatHQC}, E )$ are written in terms of the $X$-basis, then obtaining concise relations for $\Gamma \setminus \{ g_Z \}$ is somewhat tedious.
For this reason, the equational theory has been omitted.
Intuitively, this theory corresponds to writing the generalized control relations with respect to the $Z$-basis, as opposed to the $X$-basis.

\begin{figure}[t]
    \begin{framed}
    \textbf{Collection 1: Ordinary Circuit Relations}
    \vspace{-1em}
    \begin{align*}
        \scalebox{0.65}{\circimg{tikz-cache/proofs__decode_q3_lhs.pdf}}
        &\overset{\HqcSwapDef}{=}
        \scalebox{0.65}{\circimg{tikz-cache/proofs__decode_q3_mid.pdf}}
        & % -------------------------------------------------------
        \scalebox{0.65}{\circimg{tikz-cache/reductions__euler_alt_lhs_2.pdf}}
        &\overset{\text{R'(6)}}{=}
        \scalebox{0.65}{\circimg{tikz-cache/hqc__euler_rhs.pdf}}
        & % -------------------------------------------------------
        \scalebox{0.65}{\circimg{tikz-cache/hqc__eqn7_lhs.pdf}}
        &\overset{\HqcPOrd}{=}
        \scalebox{0.65}{\circimg{tikz-cache/cqc__eqn1_rhs.pdf}}
    \end{align*}
    
    \textbf{Collection 2: Generalized Control Relations}
    \vspace{-0.5em}
    \begin{align*}
        \scalebox{0.65}{\circimg{tikz-cache/hqc__eqn1_lhs.pdf}}
        &\overset{\HqcCtrlId}{=}
        \scalebox{0.65}{\circimg{tikz-cache/hqc__eqn1_rhs.pdf}}
        & % -------------------------------------------------------
        \scalebox{0.65}{\circimg{tikz-cache/hqc__eqn2_lhs.pdf}}
        &\overset{\HqcCtrlComp}{=}
        \scalebox{0.65}{\circimg{tikz-cache/hqc__eqn2_rhs.pdf}}
        & % -------------------------------------------------------
        \scalebox{0.65}{\circimg{tikz-cache/hqc__eqn3_lhs.pdf}}
        &\overset{\HqcCtrlSwap}{=}
        \scalebox{0.65}{\circimg{tikz-cache/hqc__eqn3_rhs.pdf}}
    \end{align*}
    \vspace{-1.5em}
    \begin{align*}
        \scalebox{0.65}{\circimg{tikz-cache/hqc__eqn18_lhs.pdf}}
        &\overset{\HqcFreePhase}{=}
        \scalebox{0.65}{\circimg{tikz-cache/hqc__eqn18_rhs.pdf}}
        & % -------------------------------------------------------
        \scalebox{0.65}{\circimg{tikz-cache/hqc__eqn14_lhs.pdf}}
        &\overset{\HqcCtrlDag}{=}
        \scalebox{0.65}{\circimg{tikz-cache/hqc__eqn14_rhs.pdf}}
        & % -------------------------------------------------------
        \scalebox{0.65}{\circimg{tikz-cache/hqc__eqn15_lhs_alt.pdf}}
        &\overset{\HqcCtrlPair}{=}
        \scalebox{0.65}{\circimg{tikz-cache/hqc__eqn15_rhs.pdf}}
    \end{align*}
    \vspace{-1.5em}
    \begin{align*}
        \scalebox{0.65}{\circimg{tikz-cache/hqc__eqn22_lhs.pdf}}
        &\overset{\HqcCtrlExpand}{=}
        \scalebox{0.65}{\circimg{tikz-cache/hqc__eqn22_rhs.pdf}}
        & % -------------------------------------------------------
        \scalebox{0.65}{\circimg{tikz-cache/hqc__eqn23_lhs.pdf}}
        &\overset{\HqcCtrlMerge}{=}
        \scalebox{0.65}{\circimg{tikz-cache/hqc__eqn23_rhs.pdf}}
    \end{align*}
    
    \textbf{Collection 3: Generalized Exponential Relations}
    \vspace{-1em}
    \begin{align*}
        \scalebox{0.65}{\circimg{tikz-cache/hqc__eqn9_lhs.pdf}}
        &\overset{\HqcPowZero}{=}
        \scalebox{0.65}{\circimg{tikz-cache/hqc__eqn9_rhs.pdf}}
        & % -------------------------------------------------------
        \scalebox{0.65}{\circimg{tikz-cache/hqc__eqn4_lhs.pdf}}
        &\overset{\HqcPowId}{=}
        \scalebox{0.65}{\circimg{tikz-cache/hqc__eqn4_rhs.pdf}}
        & % -------------------------------------------------------
        \scalebox{0.65}{\circimg{tikz-cache/hqc__eqn5_lhs.pdf}}
        &\overset{\HqcPowAdd}{=}
        \scalebox{0.65}{\circimg{tikz-cache/hqc__eqn5_rhs.pdf}}
        & % -------------------------------------------------------
        \scalebox{0.65}{\circimg{tikz-cache/hqc__eqn24_lhs.pdf}}
        &\overset{\HqcPowIdem}{=}
        \scalebox{0.65}{\circimg{tikz-cache/hqc__eqn24_rhs.pdf}}
    \end{align*}
    \vspace{-2.5em}
    \begin{align*}
        \scalebox{0.65}{\circimg{tikz-cache/hqc__eqn19_lhs.pdf}}
        &\overset{\HqcPowDag}{=}
        \scalebox{0.65}{\circimg{tikz-cache/hqc__eqn19_rhs.pdf}}
        & % -------------------------------------------------------
        \scalebox{0.65}{\circimg{tikz-cache/hqc__eqn21_lhs.pdf}}
        &\overset{\HqcPowExpand}{=}
        \scalebox{0.65}{\circimg{tikz-cache/hqc__eqn21_rhs.pdf}}
    \end{align*}
    
    \textbf{Collection 4: Congruence Relations (Inductively Defined)}
    \vspace{-1em}
    \begin{equation*}
        \scalebox{0.65}{\circimg{tikz-cache/hqc__equiv_lhs.pdf}}
        \approx_E
        \scalebox{0.65}{\circimg{tikz-cache/hqc__equiv_rhs.pdf}}
        %
        \qquad\Rightarrow\qquad
        %
        \scalebox{0.65}{\circimg{tikz-cache/hqc__eqn20_lhs.pdf}}
        \overset{\HqcPowCong}{=}
        \scalebox{0.65}{\circimg{tikz-cache/hqc__eqn20_rhs.pdf}}
        %
        \qquad\text{and}\qquad
        %
        \scalebox{0.65}{\circimg{tikz-cache/hqc__eqn13_lhs.pdf}}
        \overset{\HQCCtrlCong}{=}
        \scalebox{0.65}{\circimg{tikz-cache/hqc__eqn13_rhs.pdf}}
    \end{equation*}
    
    \vspace{-1.5em}
\end{framed}

    \caption{An alternative equational theory for $\CatHQC$, using only the generators in $\Gamma_H$.}
    \label{Fig:RHRelations}
\end{figure}

\vspace{1em}
\noindent
\textbf{Eliminating the Hadamard Gate}.
%
Let $g_H = \scalebox{0.65}{\circimg{tikz-cache/dagger__term5.pdf}}$.
To eliminate $g_H$, it is necessary to first introduce \cref{Eqn:HadDef} as a relation \text{R'(1)} via a Tietze transformation.
Since $g_H$ appears as the left-hand side of \text{R'(1)}, and does not appear on the right-hand side of \text{R'(1)}, then $g_H$ and \text{R'(1)} can be eliminated via a Tietze transformation. 
Unfortunately, $g_H$ still appears in relations $\{ \HqcTargDef, \HqcEuler \}$, meaning that some care must be taken before $g_H$ can be fully removed from $( \Gamma, R )$.
In turns out that the relation \HqcTargDef is redundant, as illustrated by the following derivation.
\begin{align*}
    \scalebox{0.65}{\circimg{tikz-cache/reductions__targdef_1.pdf}}
    \overset{\text{R'(1)}}{=}&
    \scalebox{0.65}{\circimg{tikz-cache/reductions__targdef_2.pdf}}
    \overset{\HqcPowZero}{=}
    \scalebox{0.65}{\circimg{tikz-cache/reductions__targdef_3.pdf}}
    \overset{\HqcPowAdd}{=}
    \scalebox{0.65}{\circimg{tikz-cache/reductions__targdef_4.pdf}}
    \\
    \overset{\HqcEuler}{=}&
    \scalebox{0.65}{\circimg{tikz-cache/reductions__targdef_5.pdf}}
    \overset{\HqcPowZero}{=}
    \scalebox{0.65}{\circimg{tikz-cache/reductions__targdef_6.pdf}}
    \overset{\HqcPowAdd}{=}
    \scalebox{0.65}{\circimg{tikz-cache/reductions__targdef_7.pdf}}
    \overset{\HqcPowAdd}{=}
    \scalebox{0.65}{\circimg{tikz-cache/reductions__targdef_8.pdf}}
    \\
    \overset{\HqcPOrd}{=}&
    \scalebox{0.65}{\circimg{tikz-cache/reductions__targdef_9.pdf}}
    \overset{\HqcPowId}{=}
    \scalebox{0.65}{\circimg{tikz-cache/reductions__targdef_10.pdf}}
    \overset{\HqcFreePhase}{=}
    \scalebox{0.65}{\circimg{tikz-cache/reductions__targdef_11.pdf}}
    \overunderset{\HQCCtrlCong}{\HqcPOrd}{=}
    \scalebox{0.65}{\circimg{tikz-cache/reductions__targdef_12.pdf}}
    \overset{\HqcCtrlId}{=}
    \scalebox{0.65}{\circimg{tikz-cache/reductions__targdef_13.pdf}}
\end{align*}
To replace \HqcEuler, the following derivation suffices.
\begin{equation*}
    \scalebox{0.65}{\circimg{tikz-cache/hqc__euler_lhs.pdf}}
    \overset{\text{R'(1)}}{=}
    \scalebox{0.65}{\circimg{tikz-cache/reductions__euler_alt_lhs_1.pdf}}
    \overset{\HqcPowAdd}{=}
    \scalebox{0.65}{\circimg{tikz-cache/reductions__euler_alt_lhs_2.pdf}}
\end{equation*}
Using this derivation, it is then possible to introduce a new relation \text{R'(6)} as in \cref{Fig:RHRelations}, which is then used to eliminate \HqcEuler.
Since \HqcTargDef has been eliminated and $g_H$ does not appear in any instance of \text{R'(6)}, then $g_H$ can be removed using \text{R'(1)} without any further modifications to the relation set.
This yields a new presentation, $( \Gamma_H, R_H )$ where $\Gamma_H$ and $R_H$ are defined as follows.
\begin{align*}
    \Gamma &= \{
        \scalebox{0.65}{\circimg{tikz-cache/dagger__term1.pdf}},
        \scalebox{0.65}{\circimg{tikz-cache/dagger__term3.pdf}},
        \;\scalebox{0.65}{\circimg{tikz-cache/dagger__term8.pdf}}
    \}
    &
    R_H &= R \setminus \{
        \HqcTargDef,
        \HqcEuler
    \} \cup \{
        \text{E'(6)}
    \}
\end{align*}
The relations in $Q_H$ are depicted in \cref{Fig:RHRelations}.
